\begin{table}[!htbp]
\centering\small
\caption{Moc testu Wilcoxona przy 12 i~24 zbiorach. Podano najmniejszą osiągalną dwustronną wartość $p$ oraz największą sumę rang $T$ zbiorów o~przeciwnym znaku różnicy, przy której wynik jest jeszcze istotny w~najostrzejszym kroku procedury Holma (14 testów, próg $0.05/14$) i~bez korekty (próg $0.05$).}
\label{tab:moc}
\begin{tabular}{rrrr}
\toprule
\textbf{zbiory} & \textbf{min.\ $p$} & \textbf{maks.\ $T$ (Holm, krok 1)} & \textbf{maks.\ $T$ (bez korekty)} \\
\midrule
12 & 0.00049 & 4 & 13 \\
24 & 1.2e-07 & 51 & 81 \\
\bottomrule
\end{tabular}
\end{table}
